\section{Disparate Impact and Error Rates}
\label{sec:note-meas-disp}

We start by reinterpreting the  ``80\% rule'' in terms of more standard statistical measures of quality of a classifier. This presents notational challenges. The terminology of ``right'' and ``wrong'', ``positive'' and ``negative'' that is used in classification is an awkward fit when dealing with majority and minority classes, and selection decisions. For \emph{notational convenience only}, we will use the convention that the protected class $X$ takes on two values: $X =0$ for the ``minority'' class and $X = 1$ for the ``default'' class. For example, in most gender-discrimination scenarios the value $0$ would be assigned to ``female'' and $1$ to ``male''. We will denote a successful binary classification outcome $C$ (say, a hiring decision) by $C =\,$\textsc{yes} and a failure by $C=\, $\textsc{no}. Finally, we will map the majority class to ``positive'' examples and the minority class to ``negative'' examples with respect to the classification outcome, all the while reminding the reader that this is merely a convenience to do the mapping, and does not reflect any judgments about the classes. The advantage of this mapping is that it renders our results more intuitive: a classifier with high ``error'' will also be one that is least biased, because it is unable to distinguish the two classes. 



Table \ref{tab:confusion} describes the \emph{confusion matrix} for a classification with respect to the above attributes where each entry is the probability of that particular pair of outcomes for data sampled from the input distribution (we use the empirical distribution when referring to a specific data set).
\begin{table}[htbp]
  \centering
  \begin{tabular}{c|c|c|}
  Outcome & $X = 0$ & $X = 1$ \\ \hline
  $C=\,$\textsc{no} & $a$ & $b$\\ \hline
$C=\,$\textsc{yes} & $c$& $d$\\ \hline
\end{tabular}
\caption{A confusion matrix\label{tab:confusion}}
\end{table}

The $80\%$ rule can then be quantified as: 
\[ \frac{c/(a + c)}{d/(b + d)} \ge 0.8 \]
Note that the traditional notion of ``accuracy'' includes terms in the numerator from both columns, and so cannot be directly compared to the $80\%$ rule. Still, other class-sensitive error metrics are known, and more directly relate to the $80\%$ rule:

\begin{defn}[{\small Class-conditioned error metrics}]
  The \emph{sensitivity} of a test (informally, its true positive rate) is defined as the conditional probability of returning \textsc{yes} on ``positive'' examples (a.k.a. the majority class). In other words, 
\[ \sens = \frac{d}{b+d} \]
  The \emph{specificity} of a test (its true negative rate) is defined as the conditional probability of returning \textsc{no} on ``negative'' examples (a.k.a. the minority) class. I.e.,
\[ \specf = \frac{a}{a+c} \]
\end{defn}
%Note that sensitivity is the same as the information-retrieval notion of \emph{recall}: for example, how many of the correct search results are actually returned by a query. 

\begin{defn}[{\small Likelihood ratio (positive)}]
The likelihood ratio positive, denoted by \lrplus, is given by 
\[ \lrplus(C, X) = \frac{\sens}{1 - \specf} = \frac{d/(b+d)}{c/(a+c)} \]
\end{defn}

We can now restate the $80\%$ rule in terms of a data set.
\begin{defn}[Disparate Impact]
%A data set admits (potential) disparate impact if there exists a mapping $g: Y \rightarrow C$ such that 
%\[ \lrplus(g(y), X) > \frac{1}{\tau} = 1.25 \]
A data set has \emph{disparate impact} if
\[ \lrplus(C, X) > \frac{1}{\tau} = 1.25 \]
\end{defn}

It will be convenient to work with the reciprocal of $\lrplus$, which we denote by 
\[ \di = \frac{1}{\lrplus(C, X)} \mbox{ .}\]
This will allow us to discuss the value associated with disparate impact before the threshold is applied.

% \mypara{Likelihood Ratio And Balanced Error Rate.}
% The \emph{balanced error rate} is the average error rate on the two classes in a binary classification problem. Using the confusion matrix in \autoref{tab:confusion}, we can write the balanced ``error'' rate (again reminding the reader that we are interpreting ``error'' as a \textsc{yes} outcome for the minority class) as
% \[ \ber(C, X)  = \frac{1}{2}\paren{\frac{b}{b+d} + \frac{c}{a+c}}\]
% For convenience, let us set $\alpha \triangleq \frac{b}{b+d}$ and $\beta \triangleq \frac{c}{c+a}$ in what follows. Then we can write 
% \[ \ber(C, X) = \frac{1}{2}(\alpha + \beta)\]
% and 
% \[ \lrplus(C, X) = \frac{1-\alpha}{\beta} \]

% Note that $\lrplus$ is a number that increases as the error decreases, the opposite of \ber.

\looseness-1\vspace*{0.1in}\mypara{Multiple classes.} Disparate impact is defined only for two classes. In general, one might imagine a multivalued class attribute (for example, like ethnicity). In this paper, we will assume that a multivalued class attribute has one value designated as the ``default'' or majority class, and will compare each of the other values pairwise to this default class. While this ignores zero-sum effects between the different class values, it reflects the current binary nature of legal thought on discrimination. A more general treatment of joint discrimination among multiple classes is beyond the scope of this work.


%%% Local Variables:
%%% mode: latex
%%% TeX-PDF-mode: t
%%% TeX-master: "bias"
%%% End:
